\begin{frame}[t]{\ev{\tagO}Existing runtimes reuse state for cloning, evaluation, and recovery.}
\vspace{0pt}
{\small
\begin{tabularx}{\textwidth}{@{}L{3.5cm}Y@{}}
\toprule
\textbf{System} & \textbf{Reused state and operation}\\
\midrule
SnowFlock & \texttt{VM\_fork(N)} clones reached VM state.\\
Kimi AgentENV & Checkpoint, resume, and branch the sandbox for reward evaluation.\\
DeepSeek DSec & Retain the sandbox and replay cached command results on resumption.\\
\bottomrule
\end{tabularx}\par}
\vspace{.2cm}
\lead{Kimi reports \textbf{51.2 million sandboxes} across training and evaluation.}
{\small A sandbox branch for judging does not imply a forked judge.\par}
\sources{SnowFlock, EuroSys 2009; Kimi K3, \S5.3.2; DeepSeek-V4, \S5.2.5.}
\qadetail{SnowFlock: Lagar-Cavilla et al., Rapid Virtual Machine Cloning for Cloud Computing, EuroSys 2009, built on Xen. Kimi K3 reports 51,219,741 sandboxes across 1,505,678 images in training and evaluation, up to 98\% of sandbox lifetime during model waits, and checkpoint/resume times as low as 133/49 ms. These are attributed infrastructure figures, not our measurements. Forking is offered for reward judging without side effects. DSec retains sandboxes when training tasks are preempted and replays cached command results on resumption. DeepSeek-V4, arXiv:2606.19348, \S5.2.5.}
\note{[0:45] SnowFlock exposed VM fork in 2009, on top of Xen. Kimi's AgentENV supports checkpointing, resumption, and sandbox branches for reward evaluation. Its report counts fifty-one million sandboxes across training and evaluation. DeepSeek's DSec retains environments and replays cached command results when a training task resumes. These operations preserve different surfaces and serve different purposes. A sandbox branch does not establish independent judgments. We ask what contract an agent workflow needs around these mechanisms. First consider why an additional execution can help.}
\end{frame}
